\documentclass[11pt,a4paper]{article}
\usepackage[utf8]{inputenc}
\usepackage[italian]{babel}
\usepackage{geometry}
\usepackage{hyperref}
\usepackage{xcolor}
\usepackage{listings}
\usepackage{amsmath}
\usepackage{graphicx}
\usepackage{titlesec}
\usepackage{fancyhdr}
\usepackage{enumitem}
\usepackage{booktabs}
\usepackage{float}

% Configurazione geometria pagina
\geometry{
    left=2.5cm,
    right=2.5cm,
    top=3cm,
    bottom=3cm
}

% Configurazione hyperref
\hypersetup{
    colorlinks=true,
    linkcolor=blue,
    filecolor=magenta,      
    urlcolor=cyan,
    pdftitle={Sistema di Autenticazione - Documentazione Tecnica},
    pdfauthor={ExtraTracker Team}
}

% Configurazione listings per codice
\lstset{
    language=JavaScript,
    basicstyle=\ttfamily\small,
    keywordstyle=\color{blue}\bfseries,
    commentstyle=\color{green!60!black},
    stringstyle=\color{red},
    numbers=left,
    numberstyle=\tiny\color{gray},
    stepnumber=1,
    numbersep=5pt,
    backgroundcolor=\color{gray!10},
    frame=single,
    breaklines=true,
    breakatwhitespace=true,
    tabsize=2,
    showstringspaces=false,
    escapeinside={\%*}{*)}
}

% Configurazione titoli
\titleformat{\section}
{\Large\bfseries\color{blue!80!black}}
{}
{0em}
{}[\titlerule]

\titleformat{\subsection}
{\large\bfseries\color{blue!60!black}}
{}
{0em}
{}

% Header e footer
\pagestyle{fancy}
\fancyhf{}
\fancyhead[L]{\leftmark}
\fancyhead[R]{\thepage}
\fancyfoot[C]{Sistema di Autenticazione - Documentazione Tecnica}

% Comandi personalizzati
\newcommand{\code}[1]{\texttt{#1}}
\newcommand{\tech}[1]{\textbf{\textit{#1}}}
\newcommand{\important}[1]{\textcolor{red!80!black}{\textbf{#1}}}

\title{
    \vspace{-2cm}
    \textbf{\Huge Sistema di Autenticazione}\\
    \Large Documentazione Tecnica Completa\\
    \large Backend e Frontend
    \vspace{0.5cm}
}

\author{ExtraTracker Development Team}
\date{\today}

\begin{document}

\maketitle
\thispagestyle{empty}

\newpage
\tableofcontents
\newpage

\section{Introduzione}

Questo documento fornisce una documentazione esaustiva del sistema di autenticazione implementato in ExtraTracker, analizzando in dettaglio tutte le tecniche utilizzate, le scelte progettuali e le implicazioni di sicurezza.

\subsection{Obiettivi del Sistema}

Il sistema di autenticazione è progettato per garantire:

\begin{itemize}
    \item \textbf{Sicurezza}: Protezione contro attacchi comuni (brute force, XSS, CSRF, token theft)
    \item \textbf{Usabilità}: Esperienza utente fluida con refresh automatico dei token
    \item \textbf{Conformità}: Rispetto dei principi GDPR e best practices di sicurezza
    \item \textbf{Scalabilità}: Architettura modulare e testabile
\end{itemize}

\subsection{Architettura Generale}

Il sistema segue un'architettura a \tech{tre livelli}:

\begin{enumerate}
    \item \textbf{Controller Layer}: Gestisce HTTP request/response, validazione input, gestione cookie
    \item \textbf{Service Layer}: Contiene la logica di business pura (hashing, generazione token, validazione)
    \item \textbf{Model Layer}: Schema database e metodi di accesso ai dati
\end{enumerate}

Questa separazione garantisce:
\begin{itemize}
    \item \textbf{Testabilità}: Ogni layer può essere testato indipendentemente
    \item \textbf{Manutenibilità}: Modifiche isolate per responsabilità
    \item \textbf{Riutilizzabilità}: I service possono essere usati da CLI, worker, etc.
\end{itemize}

\section{Backend: Architettura e Componenti}

\subsection{Controller Layer}

Il controller (\code{authController.js}) gestisce esclusivamente la comunicazione HTTP. \important{Non contiene logica di business}.

\subsubsection{Gestione Cookie Sicuri}

I token vengono memorizzati in \tech{HttpOnly cookies} per prevenire accesso da JavaScript (protezione XSS):

\begin{lstlisting}[caption=Impostazione Cookie Sicuri]
const setAuthCookies = (res, accessToken, refreshToken) => {
    // Access token cookie
    res.cookie(
        securityConfig.cookie.name,
        accessToken,
        securityConfig.cookie.options  // httpOnly: true, secure: true, sameSite: 'none'
    );

    // Refresh token cookie (path specifico!)
    res.cookie(
        securityConfig.cookie.refreshName,
        refreshToken,
        securityConfig.cookie.refreshOptions  // path: '/api/auth/refresh'
    );
};
\end{lstlisting}

\textbf{Perché HttpOnly?}
\begin{itemize}
    \item Previene accesso da \code{document.cookie} (XSS attacks)
    \item Il browser invia automaticamente i cookie nelle richieste
    \item Non accessibile da JavaScript malizioso
\end{itemize}

\textbf{Perché SameSite: 'none'?}
\begin{itemize}
    \item Necessario per richieste cross-origin (frontend e backend su domini diversi)
    \item Richiede \code{secure: true} (HTTPS in produzione)
    \item Previene attacchi CSRF quando configurato correttamente
\end{itemize}

\subsubsection{Flusso di Registrazione}

\begin{lstlisting}[caption=Registrazione Utente]
const register = asyncHandler(async (req, res) => {
    // 1. Validazione input (Zod middleware)
    // 2. Chiamata service per logica business
    const { user, accessToken, refreshToken } = 
        await authService.register(req.body);

    // 3. Genera token verifica email (non blocca risposta)
    const verificationToken = generateToken();
    const hashedToken = hashToken(verificationToken);
    
    // 4. Salva token hashato nel DB
    await User.findByIdAndUpdate(user._id, {
        emailVerificationToken: hashedToken,
        emailVerificationExpires: new Date(Date.now() + 24h),
    });

    // 5. Invia email (async, non blocca)
    emailService.sendVerificationEmail(user.email, verificationToken)
        .catch(err => console.error('Errore invio email:', err));

    // 6. Imposta cookie e rispondi
    setAuthCookies(res, accessToken, refreshToken);
    res.status(201).json({ success: true, data: { user } });
});
\end{lstlisting}

\textbf{Scelte progettuali:}
\begin{itemize}
    \item \textbf{Token hashato nel DB}: Il token di verifica email viene hashato prima di essere salvato (SHA-256). Solo il token in chiaro viene inviato via email.
    \item \textbf{Email asincrona}: L'invio email non blocca la risposta HTTP, migliorando la performance percepita.
    \item \textbf{TTL 24h}: Il token di verifica scade dopo 24 ore per bilanciare sicurezza e usabilità.
\end{itemize}

\subsubsection{Flusso di Login}

\begin{lstlisting}[caption=Login con Protezione Brute Force]
const login = asyncHandler(async (req, res) => {
    const { user, accessToken, refreshToken } = 
        await authService.login(req.body);
    
    setAuthCookies(res, accessToken, refreshToken);
    res.status(200).json({ success: true, data: { user } });
});
\end{lstlisting}

Il controller è semplice perché la logica complessa (verifica password, lockout account) è nel service layer.

\subsubsection{Refresh Token con Rotazione}

\begin{lstlisting}[caption=Refresh Token con Rotazione]
const refresh = asyncHandler(async (req, res) => {
    const refreshToken = req.cookies?.[securityConfig.cookie.refreshName];

    if (!refreshToken) {
        return res.status(401).json({
            success: false,
            error: { message: 'Sessione scaduta', code: 'NO_REFRESH_TOKEN' }
        });
    }

    // Service genera NUOVI token (rotation)
    const { accessToken, refreshToken: newRefreshToken } = 
        await authService.refreshAccessToken(refreshToken);

    // Ruota i token (il vecchio refresh token non è più valido)
    setAuthCookies(res, accessToken, newRefreshToken);

    res.status(200).json({ success: true, message: 'Token rinnovato' });
});
\end{lstlisting}

\textbf{Perché Refresh Token Rotation?}
\begin{itemize}
    \item \textbf{Sicurezza}: Se un refresh token viene rubato, può essere usato solo una volta
    \item \textbf{Rilevamento furto}: Se il vecchio token viene riutilizzato, è segno di compromissione
    \item \textbf{Invalidazione automatica}: Il vecchio token diventa immediatamente inutile
\end{itemize}

\subsection{Service Layer}

Il service layer (\code{authService.js}) contiene la \important{logica pura di business}, senza dipendenze da Express o HTTP.

\subsubsection{Password Hashing con Argon2id}

\begin{lstlisting}[caption=Hashing Password con Argon2id]
async hashPassword(password) {
    return argon2.hash(password, {
        type: argon2.argon2id,        // Tipo più sicuro
        memoryCost: 65536,            // 64MB (resiste a GPU attacks)
        timeCost: 3,                  // Iterazioni
        parallelism: 4,               // Usa 4 core CPU
        hashLength: 32,               // 32 bytes hash
    });
}
\end{lstlisting}

\textbf{Perché Argon2id invece di bcrypt o scrypt?}

\begin{table}[H]
\centering
\begin{tabular}{@{}lccc@{}}
\toprule
\textbf{Algoritmo} & \textbf{GPU Resistance} & \textbf{Side-Channel} & \textbf{Standard} \\
\midrule
bcrypt & Media & Vulnerabile & No \\
scrypt & Alta & Vulnerabile & No \\
Argon2id & Alta & Resistente & \checkmark (PHC 2015) \\
\bottomrule
\end{tabular}
\caption{Confronto algoritmi di hashing}
\end{table}

\textbf{Vantaggi di Argon2id:}
\begin{itemize}
    \item \textbf{Memory-hard}: Richiede grandi quantità di memoria, rendendo attacchi GPU costosi
    \item \textbf{Resistente a side-channel}: Protegge contro timing attacks
    \item \textbf{Vincitore Password Hashing Competition 2015}: Standard de facto
    \item \textbf{Configurabile}: Parametri ottimizzabili per performance/sicurezza
\end{itemize}

\textbf{Parametri scelti:}
\begin{itemize}
    \item \code{memoryCost: 65536} (64MB): Bilanciamento tra sicurezza e performance
    \item \code{timeCost: 3}: ~200ms per hash (accettabile per UX)
    \item \code{parallelism: 4}: Sfrutta multi-core moderni
\end{itemize}

\subsubsection{Generazione JWT Token}

\begin{lstlisting}[caption=Generazione Access Token]
generateAccessToken(user) {
    const payload = {
        sub: user._id.toString(),  // Subject: ID utente
        email: user.email,
        type: 'access',             // Distingue access da refresh
    };

    return jwt.sign(payload, securityConfig.jwt.secret, {
        expiresIn: '15m',           // Breve durata per sicurezza
        algorithm: 'HS256',         // Algoritmo simmetrico
    });
}
\end{lstlisting}

\textbf{Perché 15 minuti per access token?}
\begin{itemize}
    \item \textbf{Minimizza finestra di compromissione}: Se rubato, è valido solo 15 minuti
    \item \textbf{Refresh automatico}: Il frontend rinnova automaticamente quando scade
    \item \textbf{Bilanciamento}: Non troppo breve (troppe richieste) né troppo lungo (rischio sicurezza)
\end{itemize}

\begin{lstlisting}[caption=Generazione Refresh Token con Hash]
async generateRefreshToken(user) {
    // Parte random per unicità
    const randomPart = crypto.randomBytes(32).toString('hex');

    const payload = {
        sub: user._id.toString(),
        type: 'refresh',
        jti: randomPart,  // JWT ID unico
    };

    const token = jwt.sign(payload, securityConfig.jwt.secret, {
        expiresIn: '7d',  // Durata più lunga
        algorithm: 'HS256',
    });

    // Hash del token per salvarlo nel DB (NON salvare token in chiaro!)
    const hash = crypto
        .createHash('sha256')
        .update(token)
        .digest('hex');

    return { token, hash };
}
\end{lstlisting}

\textbf{Perché hashare il refresh token nel DB?}
\begin{itemize}
    \item \textbf{Protezione database}: Se il DB viene compromesso, i token non sono utilizzabili
    \item \textbf{Verifica integrità}: Confronto hash per validare token durante refresh
    \item \textbf{Best practice}: Mai salvare token in chiaro nel database
\end{itemize}

\subsubsection{Refresh Token con Validazione Hash}

\begin{lstlisting}[caption=Refresh Token con Validazione]
async refreshAccessToken(refreshToken) {
    // 1. Verifica firma JWT
    const payload = this.verifyToken(refreshToken, 'refresh');

    // 2. Trova utente
    const user = await User.findByRefreshToken(payload.sub);
    if (!user) {
        throw AppError.unauthorized('Sessione non valida');
    }

    // 3. Verifica hash del refresh token
    const tokenHash = crypto
        .createHash('sha256')
        .update(refreshToken)
        .digest('hex');

    if (user.refreshTokenHash !== tokenHash) {
        // Possibile furto token! Invalida tutti i refresh token
        user.refreshTokenHash = undefined;
        await user.save();
        throw AppError.unauthorized('Sessione invalidata per sicurezza');
    }

    // 4. Genera nuovi token (rotation)
    const newAccessToken = this.generateAccessToken(user);
    const { token: newRefreshToken, hash: newRefreshTokenHash } = 
        await this.generateRefreshToken(user);

    // 5. Aggiorna hash nel DB
    user.refreshTokenHash = newRefreshTokenHash;
    await user.save();

    return { user, accessToken: newAccessToken, refreshToken: newRefreshToken };
}
\end{lstlisting}

\textbf{Meccanismo di sicurezza:}
\begin{enumerate}
    \item Verifica firma JWT (garantisce autenticità)
    \item Verifica hash nel DB (garantisce che il token non sia stato revocato)
    \item Se hash non corrisponde: possibile furto, invalida immediatamente
    \item Genera nuovi token (rotation) per invalidare il vecchio
\end{enumerate}

\subsubsection{Login con Protezione Brute Force}

\begin{lstlisting}[caption=Login con Account Lockout]
async login(data) {
    const { email, password } = data;

    // Trova utente con campi nascosti (password, failedAttempts)
    const user = await User.findForLogin(email);

    if (!user) {
        // Messaggio generico per non rivelare se email esiste
        throw AppError.unauthorized('Email o password non corretti');
    }

    // Verifica se account bloccato
    if (user.isLocked) {
        const minutesLeft = Math.ceil((user.lockUntil - Date.now()) / 60000);
        throw AppError.unauthorized(
            `Account temporaneamente bloccato. Riprova tra ${minutesLeft} minuti`
        );
    }

    // Verifica password
    const isValidPassword = await this.verifyPassword(user.password, password);

    if (!isValidPassword) {
        // Incrementa tentativi falliti
        await user.incrementFailedAttempts();
        throw AppError.unauthorized('Email o password non corretti');
    }

    // Login riuscito: resetta contatore
    await user.resetFailedAttempts();

    // Genera token
    const accessToken = this.generateAccessToken(user);
    const { token: refreshToken, hash: refreshTokenHash } = 
        await this.generateRefreshToken(user);

    user.refreshTokenHash = refreshTokenHash;
    await user.save();

    return { user, accessToken, refreshToken };
}
\end{lstlisting}

\textbf{Protezioni implementate:}
\begin{itemize}
    \item \textbf{Account Lockout}: Dopo 5 tentativi falliti, account bloccato per 15 minuti
    \item \textbf{Messaggi generici}: Non rivela se l'email esiste (previene enumeration attacks)
    \item \textbf{Reset automatico}: Dopo login riuscito, contatore resettato
\end{itemize}

\subsection{Middleware di Autenticazione}

\subsubsection{requireAuth Middleware}

\begin{lstlisting}[caption=Middleware Require Auth]
const requireAuth = async (req, res, next) => {
    try {
        // 1. Estrai token dal cookie HttpOnly
        const token = req.cookies?.[securityConfig.cookie.name];

        if (!token) {
            throw AppError.unauthorized('Accesso negato. Effettua il login.');
        }

        // 2. Verifica token
        const payload = authService.verifyToken(token, 'access');

        // 3. Aggiungi dati utente alla request
        req.user = {
            id: payload.sub,
            email: payload.email,
        };

        next();
    } catch (error) {
        next(error);  // Passa all'error handler
    }
};
\end{lstlisting}

\textbf{Caratteristiche:}
\begin{itemize}
    \item \textbf{Estrazione da cookie}: Non richiede header Authorization (più sicuro)
    \item \textbf{Verifica tipo token}: Assicura che sia un access token, non refresh
    \item \textbf{Iniezione req.user}: Dati utente disponibili in tutti i controller successivi
\end{itemize}

\subsubsection{optionalAuth Middleware}

\begin{lstlisting}[caption=Middleware Optional Auth]
const optionalAuth = async (req, res, next) => {
    try {
        const token = req.cookies?.[securityConfig.cookie.name];

        if (token) {
            const payload = authService.verifyToken(token, 'access');
            req.user = {
                id: payload.sub,
                email: payload.email,
            };
        }
    } catch {
        // Ignora errori, l'utente semplicemente non è autenticato
        req.user = null;
    }

    next();
};
\end{lstlisting}

\textbf{Use case:}
\begin{itemize}
    \item Route che funzionano sia per utenti autenticati che non
    \item Logout: deve funzionare anche se il token è scaduto
\end{itemize}

\subsection{Validazione Input con Zod}

\subsubsection{Schema di Validazione}

\begin{lstlisting}[caption=Schema Zod per Registrazione]
const emailSchema = z
    .string({
        required_error: 'Email obbligatoria',
        invalid_type_error: 'Email deve essere una stringa',
    })
    .email('Formato email non valido')
    .min(5, 'Email troppo corta')
    .max(255, 'Email troppo lunga')
    .toLowerCase()  // Normalizza a minuscolo
    .trim();        // Rimuovi spazi

const passwordSchema = z
    .string({
        required_error: 'Password obbligatoria',
    })
    .min(8, 'Password deve essere almeno 8 caratteri')
    .max(128, 'Password troppo lunga')
    .regex(
        /^(?=.*[a-z])(?=.*[A-Z])(?=.*\d)/,
        'Password deve contenere almeno: 1 maiuscola, 1 minuscola, 1 numero'
    );

const registerSchema = z.object({
    email: emailSchema,
    password: passwordSchema,
    acceptTerms: z
        .boolean({
            required_error: 'Devi accettare i termini e condizioni',
        })
        .refine((val) => val === true, {
            message: 'Devi accettare i termini e condizioni',
        }),
});
\end{lstlisting}

\textbf{Perché Zod invece di validazione manuale?}
\begin{itemize}
    \item \textbf{Type-safe}: Genera tipi TypeScript automaticamente
    \item \textbf{Composabile}: Schemi riutilizzabili e estendibili
    \item \textbf{Messaggi chiari}: Errori strutturati per il frontend
    \item \textbf{Previene NoSQL Injection}: Valida tipi rigorosi prima del DB
    \item \textbf{Sanitizzazione automatica}: \code{toLowerCase()}, \code{trim()} applicati automaticamente
\end{itemize}

\subsubsection{Middleware di Validazione}

\begin{lstlisting}[caption=Middleware Validazione Zod]
const validateMiddleware = (schema) => {
    return (req, res, next) => {
        const result = schema.safeParse(req.body);
        
        if (!result.success) {
            return res.status(400).json({
                success: false,
                error: {
                    message: 'Errore di validazione',
                    code: 'VALIDATION_ERROR',
                    details: result.error.errors.map((err) => ({
                        field: err.path.join('.'),
                        message: err.message,
                    })),
                },
            });
        }
        
        // Sostituisci req.body con dati validati e sanitizzati
        req.body = result.data;
        next();
    };
};
\end{lstlisting}

\textbf{Vantaggi:}
\begin{itemize}
    \item \textbf{Validazione precoce}: Blocca richieste malformate prima del service
    \item \textbf{Sanitizzazione}: \code{req.body} contiene solo dati validati
    \item \textbf{Errori strutturati}: Frontend può mostrare errori per campo
\end{itemize}

\subsection{Rate Limiting}

\subsubsection{Configurazione Rate Limiter}

\begin{lstlisting}[caption=Rate Limiter per Autenticazione]
const authLimiter = rateLimit({
    windowMs: 15 * 60 * 1000,  // 15 minuti
    max: 15,                    // 15 tentativi
    keyGenerator: (req) => getClientIpForRateLimit(req),
    skipSuccessfulRequests: true,  // Non conta login riusciti
    handler: (req, res) => {
        res.status(429).json({
            success: false,
            error: {
                message: 'Troppi tentativi di accesso. Account temporaneamente bloccato.',
                code: 'TOO_MANY_ATTEMPTS',
                retryAfter: 15 * 60,
            },
        });
    },
});
\end{lstlisting}

\textbf{Protezioni:}
\begin{itemize}
    \item \textbf{Brute Force}: Limita tentativi di login per IP
    \item \textbf{Enumeration}: Previene discovery di email esistenti
    \item \textbf{DDoS parziale}: Limita richieste totali per IP
\end{itemize}

\textbf{Perché 15 tentativi?}
\begin{itemize}
    \item Bilanciamento tra sicurezza e usabilità
    \item Considera utenti che dimenticano password
    \item Combinato con account lockout (5 tentativi) fornisce doppia protezione
\end{itemize}

\subsubsection{Gestione IP con Proxy}

\begin{lstlisting}[caption=Estrazione IP Reale con Proxy]
const getClientIpForRateLimit = (req) => {
    const remoteAddress = req?.socket?.remoteAddress;
    const isLoopback = 
        remoteAddress === '127.0.0.1' ||
        remoteAddress === '::1';

    const xForwardedFor = req.headers['x-forwarded-for'];
    
    // Accetta X-Forwarded-For solo se il peer diretto è loopback
    // (previene header spoofing)
    if (isLoopback && typeof xForwardedFor === 'string') {
        return xForwardedFor.split(',')[0].trim();
    }

    return req.ip;
};
\end{lstlisting}

\textbf{Sicurezza:}
\begin{itemize}
    \item \textbf{Previene header spoofing}: Accetta X-Forwarded-For solo da proxy fidati
    \item \textbf{Supporta Vite dev server}: In sviluppo, Vite fa da proxy
    \item \textbf{Produzione-ready}: Funziona con reverse proxy (Nginx, Cloudflare)
\end{itemize}

\subsection{Modello Utente}

\subsubsection{Schema Mongoose}

\begin{lstlisting}[caption=Schema Utente con Campi Sensibili Nascosti]
const userSchema = new mongoose.Schema({
    email: {
        type: String,
        required: true,
        unique: true,
        lowercase: true,
        trim: true,
        index: true,
    },
    
    password: {
        type: String,
        required: true,
        select: false,  // Non includere nelle query di default
    },
    
    refreshTokenHash: {
        type: String,
        select: false,  // Hash del refresh token
    },
    
    failedLoginAttempts: {
        type: Number,
        default: 0,
        select: false,  // Campo sensibile
    },
    
    lockUntil: {
        type: Date,
        select: false,  // Campo sensibile
    },
    
    // ... altri campi
}, {
    timestamps: true,  // createdAt, updatedAt automatici
    toJSON: {
        transform: (doc, ret) => {
            // Rimuovi campi sensibili quando converti in JSON
            delete ret.password;
            delete ret.refreshTokenHash;
            delete ret.failedLoginAttempts;
            delete ret.lockUntil;
            return ret;
        },
    },
});
\end{lstlisting}

\textbf{Protezioni:}
\begin{itemize}
    \item \textbf{select: false}: Campi sensibili non inclusi nelle query di default
    \item \textbf{toJSON transform}: Rimozione automatica campi sensibili
    \item \textbf{Metodi statici}: \code{findForLogin()} include campi necessari esplicitamente
\end{itemize}

\subsubsection{Account Lockout}

\begin{lstlisting}[caption=Metodi per Gestione Lockout]
userSchema.methods.incrementFailedAttempts = async function () {
    // Se lockout è scaduto, resetta
    if (this.lockUntil && this.lockUntil < Date.now()) {
        return this.updateOne({
            $set: { failedLoginAttempts: 1 },
            $unset: { lockUntil: 1 },
        });
    }

    const updates = { $inc: { failedLoginAttempts: 1 } };

    // Se raggiunti 5 tentativi, blocca per 15 minuti
    if (this.failedLoginAttempts + 1 >= 5) {
        updates.$set = { lockUntil: Date.now() + 15 * 60 * 1000 };
    }

    return this.updateOne(updates);
};

userSchema.methods.resetFailedAttempts = async function () {
    return this.updateOne({
        $set: { failedLoginAttempts: 0, lastLoginAt: new Date() },
        $unset: { lockUntil: 1 },
    });
};
\end{lstlisting}

\textbf{Meccanismo:}
\begin{itemize}
    \item Dopo 5 tentativi falliti, account bloccato per 15 minuti
    \item Reset automatico dopo login riuscito
    \item Reset automatico se lockout scaduto
\end{itemize}

\section{Frontend: Architettura e Componenti}

\subsection{API Client con Interceptor}

\subsubsection{Configurazione Axios}

\begin{lstlisting}[caption=Configurazione Axios con Credentials]
const axiosInstance = axios.create({
    baseURL: API_BASE_URL,
    headers: {
        'Content-Type': 'application/json',
    },
    // CRITICO: Necessario per inviare/ricevere cookies cross-origin
    withCredentials: true,
    timeout: 25000,
});
\end{lstlisting}

\textbf{Perché withCredentials: true?}
\begin{itemize}
    \item \textbf{Cookie cross-origin}: Permette invio cookie HttpOnly tra domini diversi
    \item \textbf{Sicurezza}: Il browser gestisce automaticamente i cookie
    \item \textbf{Requisito CORS}: Il backend deve avere \code{credentials: true} in CORS
\end{itemize}

\subsubsection{Interceptor per Refresh Automatico}

\begin{lstlisting}[caption=Interceptor Refresh Token Automatico]
// Flag per evitare loop infiniti
let isRefreshing = false;
let failedQueue = [];

const processQueue = (error) => {
    failedQueue.forEach((prom) => {
        if (error) {
            prom.reject(error);
        } else {
            prom.resolve();
        }
    });
    failedQueue = [];
};

axiosInstance.interceptors.response.use(
    (response) => response,
    
    async (error) => {
        const originalRequest = error.config;
        const requestUrl = originalRequest?.url || '';

        // NON fare refresh per endpoint di auth
        const shouldSkipRefresh = 
            ['/auth/login', '/auth/register', '/auth/refresh', '/auth/logout']
            .some(url => requestUrl.includes(url));
        
        if (error.response?.status === 401 && 
            !originalRequest._retry && 
            !shouldSkipRefresh) {
            
            // Se stiamo già facendo refresh, accoda la richiesta
            if (isRefreshing) {
                return new Promise((resolve, reject) => {
                    failedQueue.push({ resolve, reject });
                }).then(() => axiosInstance(originalRequest));
            }

            originalRequest._retry = true;
            isRefreshing = true;

            try {
                // Prova a fare refresh del token
                await axiosInstance.post('/auth/refresh');
                processQueue(null);
                
                // Riprova la richiesta originale
                return axiosInstance(originalRequest);
            } catch (refreshError) {
                processQueue(refreshError);
                
                // Refresh fallito: emetti evento per logout
                window.dispatchEvent(new CustomEvent('auth:sessionExpired'));
                
                return Promise.reject(refreshError);
            } finally {
                isRefreshing = false;
            }
        }

        return Promise.reject(error);
    }
);
\end{lstlisting}

\textbf{Meccanismo:}
\begin{enumerate}
    \item Richiesta riceve 401 (token scaduto)
    \item Se non è già in refresh, chiama \code{/auth/refresh}
    \item Se altre richieste arrivano durante refresh, vengono accodate
    \item Dopo refresh riuscito, tutte le richieste accodate vengono rieseguite
    \item Se refresh fallisce, emette evento per logout
\end{enumerate}

\textbf{Vantaggi:}
\begin{itemize}
    \item \textbf{Trasparente}: L'utente non nota il refresh
    \item \textbf{Efficiente}: Una sola chiamata refresh anche con multiple richieste simultanee
    \item \textbf{Resiliente}: Gestisce correttamente errori e race conditions
\end{itemize}

\subsubsection{Interceptor per Error Handling}

\begin{lstlisting}[caption=Interceptor Error Handling Centralizzato]
axiosInstance.interceptors.response.use(
    (response) => response,
    
    (error) => {
        const status = error.response?.status;
        const requestUrl = error.config?.url || '';
        
        // Skip toast per endpoint di auth (hanno gestione custom)
        const isAuthEndpoint = 
            ['/auth/login', '/auth/register', '/auth/refresh']
            .some(url => requestUrl.includes(url));
        
        // Skip toast per 401 (gestito dal refresh token flow)
        if (status === 401) {
            return Promise.reject(error);
        }

        // Estrai messaggio errore
        let errorMessage = 'Si è verificato un errore imprevisto';
        
        if (error.response?.data?.error?.message) {
            errorMessage = error.response.data.error.message;
        } else if (error.message === 'Network Error') {
            errorMessage = 'Errore di connessione. Verifica la tua rete.';
        }

        // Mostra toast di errore (solo se non è un endpoint di auth)
        if (!isAuthEndpoint) {
            emitToast.error(errorMessage, {
                title: 'Errore',
                duration: 6000,
            });
        }

        return Promise.reject(error);
    }
);
\end{lstlisting}

\textbf{Pattern: Centralized Error Handling}
\begin{itemize}
    \item \textbf{DRY}: Non devi gestire errori in ogni componente
    \item \textbf{Consistenza}: Tutti gli errori appaiono nello stesso modo
    \item \textbf{Manutenibilità}: Cambio formato toast in un solo posto
\end{itemize}

\subsection{Auth Context}

\subsubsection{Provider e Stato}

\begin{lstlisting}[caption=Auth Context Provider]
export const AuthProvider = ({ children }) => {
    const [state, setState] = useState({
        user: null,
        isAuthenticated: false,
        isLoading: true,  // Inizia con loading per verificare sessione
        error: null,
    });

    const initialCheckDoneRef = useRef(false);
    const checkInFlightRef = useRef(false);
    const lastCheckAtRef = useRef(0);

    // Verifica sessione all'avvio
    useEffect(() => {
        if (initialCheckDoneRef.current) return;
        initialCheckDoneRef.current = true;
        checkAuth();
    }, [checkAuth]);

    // Gestisci evento sessione scaduta
    useEffect(() => {
        const handleSessionExpired = () => {
            setState({
                user: null,
                isAuthenticated: false,
                isLoading: false,
                error: 'Sessione scaduta. Effettua nuovamente il login.',
            });
        };

        window.addEventListener('auth:sessionExpired', handleSessionExpired);
        return () => {
            window.removeEventListener('auth:sessionExpired', handleSessionExpired);
        };
    }, []);

    // ... metodi login, register, logout
};
\end{lstlisting}

\textbf{Caratteristiche:}
\begin{itemize}
    \item \textbf{Verifica automatica}: Controlla sessione all'avvio app
    \item \textbf{Protezione doppio mount}: \code{useRef} previene chiamate duplicate (React StrictMode)
    \item \textbf{Event-driven}: Ascolta evento \code{auth:sessionExpired} emesso da apiClient
    \item \textbf{Loading state}: Mostra loading durante verifica iniziale
\end{itemize}

\subsubsection{Metodo Login}

\begin{lstlisting}[caption=Metodo Login nel Context]
const login = useCallback(async (data) => {
    setState((prev) => ({ ...prev, isLoading: true, error: null }));

    try {
        const response = await authService.login(data);

        if (response.success && response.data?.user) {
            setState({
                user: response.data.user,
                isAuthenticated: true,
                isLoading: false,
                error: null,
            });
        } else {
            setState((prev) => ({
                ...prev,
                isLoading: false,
                error: response.error?.message || 'Errore durante il login',
            }));
        }

        return response;
    } catch (err) {
        const message = err instanceof Error ? err.message : 'Errore durante il login';
        setState((prev) => ({
            ...prev,
            isLoading: false,
            error: message,
        }));
        throw err;
    }
}, []);
\end{lstlisting}

\textbf{Pattern:}
\begin{itemize}
    \item \textbf{Optimistic UI}: Aggiorna stato solo se risposta successo
    \item \textbf{Error handling}: Cattura errori e aggiorna stato
    \item \textbf{Return response}: Permette al componente di gestire successo/errore
\end{itemize}

\subsubsection{Protected Route Component}

\begin{lstlisting}[caption=Componente Protected Route]
export const ProtectedRoute = ({ children, redirectTo = '/login' }) => {
    const { isAuthenticated, isLoading } = useAuth();
    const location = useLocation();

    // Mostra loading mentre verifica
    if (isLoading) {
        return (
            <div className="flex items-center justify-center min-h-screen">
                <div className="w-12 h-12 border-2 rounded-full animate-spin" />
            </div>
        );
    }

    // Se non autenticato, redirect a login salvando la pagina corrente
    if (!isAuthenticated) {
        return <Navigate to={redirectTo} state={{ from: location }} replace />;
    }

    return <>{children}</>;
};
\end{lstlisting}

\textbf{Caratteristiche:}
\begin{itemize}
    \item \textbf{Loading state}: Mostra spinner durante verifica
    \item \textbf{Redirect con state}: Salva pagina corrente per redirect dopo login
    \item \textbf{Replace navigation}: Non aggiunge entry alla history
\end{itemize}

\subsection{Validazione Frontend}

\subsubsection{Schemi Zod Identici al Backend}

\begin{lstlisting}[caption=Schema Zod Frontend]
export const emailSchema = z
    .string()
    .min(1, 'Email obbligatoria')
    .email('Formato email non valido')
    .max(255, 'Email troppo lunga');

export const passwordSchema = z
    .string()
    .min(8, 'Password deve essere almeno 8 caratteri')
    .max(128, 'Password troppo lunga')
    .regex(
        /^(?=.*[a-z])(?=.*[A-Z])(?=.*\d)/,
        'Deve contenere almeno: 1 maiuscola, 1 minuscola, 1 numero'
    );

export const registerSchema = z.object({
    email: emailSchema,
    password: passwordSchema,
    confirmPassword: z.string().min(1, 'Conferma password obbligatoria'),
    acceptTerms: z.boolean().refine((val) => val === true, {
        message: 'Devi accettare i termini e condizioni',
    }),
}).refine((data) => data.password === data.confirmPassword, {
    message: 'Le password non coincidono',
    path: ['confirmPassword'],
});
\end{lstlisting}

\textbf{Perché validazione anche nel frontend?}
\begin{itemize}
    \item \textbf{Feedback immediato}: L'utente vede errori prima di inviare
    \item \textbf{Riduce chiamate API}: Blocca richieste malformate
    \item \textbf{UX migliore}: Validazione in tempo reale durante digitazione
    \item \textbf{Consistenza}: Stessi schemi del backend garantiscono coerenza
\end{itemize}

\textbf{Nota importante:} La validazione frontend \important{non sostituisce} quella backend. Il backend deve sempre validare per sicurezza.

\section{Sicurezza: Tecniche e Best Practices}

\subsection{Protezione contro Attacchi Comuni}

\subsubsection{XSS (Cross-Site Scripting)}

\textbf{Protezioni implementate:}
\begin{itemize}
    \item \textbf{HttpOnly cookies}: Token non accessibili da JavaScript
    \item \textbf{Content Security Policy}: Headers CSP nel backend
    \item \textbf{Sanitizzazione input}: Zod valida e sanitizza tutti gli input
    \item \textbf{Escape output}: React escape automaticamente i valori
\end{itemize}

\subsubsection{CSRF (Cross-Site Request Forgery)}

\textbf{Protezioni implementate:}
\begin{itemize}
    \item \textbf{SameSite cookies}: \code{sameSite: 'none'} con \code{secure: true}
    \item \textbf{CORS configurato}: Solo origini permesse possono fare richieste
    \item \textbf{Double Submit Cookie}: I cookie HttpOnly non sono leggibili da JavaScript malizioso
\end{itemize}

\subsubsection{Brute Force Attacks}

\textbf{Protezioni implementate:}
\begin{enumerate}
    \item \textbf{Rate limiting}: 15 tentativi per 15 minuti per IP
    \item \textbf{Account lockout}: 5 tentativi falliti bloccano account per 15 minuti
    \item \textbf{Messaggi generici}: Non rivela se email esiste
\end{enumerate}

\subsubsection{Token Theft}

\textbf{Protezioni implementate:}
\begin{enumerate}
    \item \textbf{HttpOnly cookies}: Token non accessibili da JavaScript
    \item \textbf{Refresh token rotation}: Token rubato può essere usato solo una volta
    \item \textbf{Hash nel DB}: Refresh token hashato, non in chiaro
    \item \textbf{Validazione hash}: Verifica integrità token durante refresh
    \item \textbf{Invalidazione automatica}: Se hash non corrisponde, invalida immediatamente
\end{enumerate}

\subsubsection{Password Attacks}

\textbf{Protezioni implementate:}
\begin{enumerate}
    \item \textbf{Argon2id}: Algoritmo memory-hard resistente a GPU attacks
    \item \textbf{Parametri elevati}: 64MB memoria, 3 iterazioni, 4 parallelism
    \item \textbf{Password policy}: Minimo 8 caratteri, maiuscola, minuscola, numero
    \item \textbf{Never in plaintext}: Password mai salvate in chiaro
\end{enumerate}

\subsection{Configurazione Sicurezza}

\subsubsection{JWT Configuration}

\begin{lstlisting}[caption=Configurazione JWT]
jwt: {
    secret: process.env.JWT_SECRET || 'CHANGE_ME_IN_PRODUCTION',
    accessTokenExpiry: '15m',      // Breve durata
    refreshTokenExpiry: '7d',      // Durata più lunga
    algorithm: 'HS256',            // Algoritmo simmetrico
}
\end{lstlisting}

\textbf{Best practices:}
\begin{itemize}
    \item \textbf{Secret forte}: Minimo 32 caratteri random in produzione
    \item \textbf{Environment variable}: Secret mai committato nel codice
    \item \textbf{Access token breve}: Minimizza finestra compromissione
    \item \textbf{Refresh token lungo}: Bilanciamento sicurezza/UX
\end{itemize}

\subsubsection{Cookie Configuration}

\begin{lstlisting}[caption=Configurazione Cookie]
cookie: {
    name: 'accessToken',
    refreshName: 'refreshToken',
    
    options: {
        httpOnly: true,              // Non accessibile da JS
        secure: isProduction,        // Solo HTTPS in produzione
        sameSite: isProduction ? 'none' : 'lax',  // Cross-origin support
        path: '/',
        maxAge: 15 * 60 * 1000,      // 15 minuti
    },
    
    refreshOptions: {
        httpOnly: true,
        secure: isProduction,
        sameSite: isProduction ? 'none' : 'lax',
        path: '/api/auth/refresh',   // Path specifico
        maxAge: 7 * 24 * 60 * 60 * 1000,  // 7 giorni
    },
}
\end{lstlisting}

\textbf{Scelte:}
\begin{itemize}
    \item \textbf{Path specifico refresh}: Limita scope del cookie refresh token
    \item \textbf{SameSite: 'none'}: Necessario per cross-origin, richiede HTTPS
    \item \textbf{Secure in produzione}: Cookie inviati solo su HTTPS
\end{itemize}

\subsubsection{CORS Configuration}

\begin{lstlisting}[caption=Configurazione CORS]
cors: {
    origin: (origin, callback) => {
        const allowedOrigins = [
            process.env.FRONTEND_URL,
            'http://localhost:5173',
            'http://localhost:5174',
        ].filter(Boolean);
        
        // Permetti richieste senza origin (Postman, curl)
        if (!origin) {
            return callback(null, true);
        }
        
        // Controlla lista
        if (allowedOrigins.includes(origin)) {
            return callback(null, true);
        }
        
        callback(new Error('Not allowed by CORS'));
    },
    
    credentials: true,  // CRITICO: necessario per cookies
    allowedHeaders: ['Content-Type', 'Authorization'],
    methods: ['GET', 'POST', 'PUT', 'DELETE', 'PATCH', 'OPTIONS'],
}
\end{lstlisting}

\textbf{Importante:}
\begin{itemize}
    \item \textbf{credentials: true}: Necessario per inviare cookie cross-origin
    \item \textbf{Origin whitelist}: Solo origini permesse possono fare richieste
    \item \textbf{No wildcard}: Non usare \code{origin: '*'} con credentials
\end{itemize}

\section{Flussi di Autenticazione}

\subsection{Flusso di Registrazione}

\begin{enumerate}
    \item \textbf{Frontend}: Utente compila form, validazione Zod client-side
    \item \textbf{Frontend}: Invia \code{POST /api/auth/register} con \code{withCredentials: true}
    \item \textbf{Backend}: Rate limiter verifica limite richieste
    \item \textbf{Backend}: Validazione Zod server-side
    \item \textbf{Backend}: Service verifica email non esistente
    \item \textbf{Backend}: Service hasha password con Argon2id
    \item \textbf{Backend}: Crea utente nel database
    \item \textbf{Backend}: Genera access token (15m) e refresh token (7d)
    \item \textbf{Backend}: Hash refresh token e salva nel DB
    \item \textbf{Backend}: Genera token verifica email (hashato nel DB)
    \item \textbf{Backend}: Imposta cookie HttpOnly con token
    \item \textbf{Backend}: Invia email verifica (async, non blocca)
    \item \textbf{Backend}: Risponde con dati utente (senza password)
    \item \textbf{Frontend}: Aggiorna AuthContext con utente
    \item \textbf{Frontend}: Redirect a dashboard
\end{enumerate}

\subsection{Flusso di Login}

\begin{enumerate}
    \item \textbf{Frontend}: Utente inserisce email/password
    \item \textbf{Frontend}: Validazione Zod client-side
    \item \textbf{Frontend}: Invia \code{POST /api/auth/login}
    \item \textbf{Backend}: Rate limiter verifica limite (15 tentativi/15min)
    \item \textbf{Backend}: Validazione Zod server-side
    \item \textbf{Backend}: Service trova utente con \code{findForLogin()} (include password)
    \item \textbf{Backend}: Verifica account non bloccato (\code{isLocked})
    \item \textbf{Backend}: Verifica password con Argon2
    \item \textbf{Backend}: Se password errata:
        \begin{itemize}
            \item Incrementa \code{failedLoginAttempts}
            \item Se $\geq$ 5, imposta \code{lockUntil} (15 minuti)
            \item Lancia errore generico
        \end{itemize}
    \item \textbf{Backend}: Se password corretta:
        \begin{itemize}
            \item Resetta \code{failedLoginAttempts}
            \item Aggiorna \code{lastLoginAt}
            \item Genera nuovi token
            \item Salva hash refresh token nel DB
        \end{itemize}
    \item \textbf{Backend}: Imposta cookie HttpOnly
    \item \textbf{Backend}: Risponde con dati utente
    \item \textbf{Frontend}: Aggiorna AuthContext
    \item \textbf{Frontend}: Redirect a dashboard
\end{enumerate}

\subsection{Flusso di Refresh Token}

\begin{enumerate}
    \item \textbf{Frontend}: Richiesta API riceve 401 (token scaduto)
    \item \textbf{Frontend}: Interceptor cattura errore
    \item \textbf{Frontend}: Se non è già in refresh, chiama \code{POST /api/auth/refresh}
    \item \textbf{Frontend}: Altre richieste simultanee vengono accodate
    \item \textbf{Backend}: Estrae refresh token da cookie
    \item \textbf{Backend}: Verifica firma JWT
    \item \textbf{Backend}: Trova utente con \code{findByRefreshToken()}
    \item \textbf{Backend}: Hash refresh token ricevuto
    \item \textbf{Backend}: Confronta hash con quello nel DB
    \item \textbf{Backend}: Se hash non corrisponde:
        \begin{itemize}
            \item Possibile furto token!
            \item Invalida \code{refreshTokenHash} nel DB
            \item Lancia errore
        \end{itemize}
    \item \textbf{Backend}: Se hash corrisponde:
        \begin{itemize}
            \item Genera NUOVI token (rotation)
            \item Hash nuovo refresh token
            \item Aggiorna \code{refreshTokenHash} nel DB
            \item Imposta nuovi cookie
        \end{itemize}
    \item \textbf{Frontend}: Se refresh riuscito, riesegue richiesta originale
    \item \textbf{Frontend}: Se refresh fallito, emette evento \code{auth:sessionExpired}
    \item \textbf{Frontend}: AuthContext ascolta evento e fa logout
\end{enumerate}

\subsection{Flusso di Logout}

\begin{enumerate}
    \item \textbf{Frontend}: Utente clicca logout
    \item \textbf{Frontend}: Chiama \code{POST /api/auth/logout}
    \item \textbf{Backend}: Se utente autenticato, invalida \code{refreshTokenHash} nel DB
    \item \textbf{Backend}: Cancella cookie (access e refresh)
    \item \textbf{Backend}: Risponde successo
    \item \textbf{Frontend}: Aggiorna AuthContext (user = null, isAuthenticated = false)
    \item \textbf{Frontend}: Redirect a login
\end{enumerate}

\section{Considerazioni Finali}

\subsection{Scelte Progettuali Riepilogate}

\begin{table}[H]
\centering
\begin{tabular}{@{}ll@{}}
\toprule
\textbf{Scelta} & \textbf{Motivazione} \\
\midrule
HttpOnly Cookies & Protezione XSS, gestione automatica browser \\
Argon2id & Resistenza GPU + side-channel, standard moderno \\
JWT con rotazione & Sicurezza token, rilevamento furti \\
Refresh automatico & UX trasparente, sicurezza mantenuta \\
Zod validazione & Type-safe, composabile, previene injection \\
Rate limiting & Protezione brute force, DDoS parziale \\
Account lockout & Doppia protezione brute force \\
SameSite: 'none' & Supporto cross-origin (necessario) \\
\bottomrule
\end{tabular}
\caption{Riepilogo scelte progettuali}
\end{table}

\subsection{Best Practices Applicate}

\begin{itemize}
    \item \textbf{Defense in Depth}: Multiple layer di sicurezza
    \item \textbf{Principle of Least Privilege}: Token con durata minima necessaria
    \item \textbf{Fail Secure}: In caso di errore, sistema si chiude in modo sicuro
    \item \textbf{Never Trust Input}: Validazione sempre, frontend e backend
    \item \textbf{Secure by Default}: Configurazioni sicure di default
    \item \textbf{Separation of Concerns}: Controller/Service/Model separati
\end{itemize}

\subsection{Possibili Miglioramenti Futuri}

\begin{itemize}
    \item \textbf{2FA/MFA}: Autenticazione a due fattori
    \item \textbf{Device fingerprinting}: Rilevamento dispositivi sospetti
    \item \textbf{Session management}: Dashboard sessioni attive
    \item \textbf{IP whitelisting}: Per account critici
    \item \textbf{Biometric auth}: Supporto WebAuthn
    \item \textbf{OAuth2/SSO}: Integrazione con provider esterni
\end{itemize}

\section{Conclusioni}

Il sistema di autenticazione implementato in ExtraTracker segue best practices moderne di sicurezza, bilanciando protezione e usabilità. L'architettura modulare garantisce manutenibilità e testabilità, mentre le multiple layer di sicurezza proteggono contro attacchi comuni.

Le scelte tecniche (Argon2id, JWT con rotazione, HttpOnly cookies, rate limiting) sono state selezionate per fornire il miglior compromesso tra sicurezza, performance e esperienza utente.

\vspace{1cm}

\textit{Documento generato il \today}

\end{document}
